\section*{Chapitre \ref{ch:m2:bond-futures} --- Contrats à terme sur obligations}

\begin{solution}{exo:m2:bond-futures:1}
$111.765625 \times 0.8870 + 1.453125 = 100.589234$ pour 100~; 100\,589.23~USD par
contrat de 100\,000~USD de valeur nominale.
\end{solution}

\begin{solution}{exo:m2:bond-futures:2}
$99.132 - 111.765625 \times 0.8870 = -0.004$, soit environ $-0.14$ d'un 32e. La \omterm{def:m2:bond-futures:irr}{base brute}
compare le prix au comptant du jour au prix à terme ajusté par le \omterm{def:m2:bond-futures:cf}{facteur de concordance}~;
le future est un forward, donc l'écart contient aussi le carry jusqu'à la livraison.
Ici le carry est négatif (le coupon de l'obligation, 3.875~\%, rapporte moins qu'un repo à 3.90~\%
sur son prix), donc le forward est au-dessus du comptant et la \omterm{def:m2:bond-futures:irr}{base brute} peut
être négative alors que la \omterm{def:m2:bond-futures:irr}{base nette} est positive.
\end{solution}

\begin{solution}{exo:m2:bond-futures:3}
Du 1\up{er} décembre 2026 au 15 novembre 2033~: 6 ans 11 mois, arrondis à
6 ans 9 mois, donc $n = 6$, $z = 9$, $v = 3$. Pour l'obligation à 4~\%, le facteur vaut 0.8902. Pour l'obligation à 6~\%
vaut 0.9999, et non 1, car la formule la valorise comme si son prochain coupon
était à trois mois et sa durée arrondie par défaut~; seule une obligation à 6~\% dont la durée est un
nombre entier de trimestres à compter du premier jour du mois a un facteur
exactement égal à 1.
\end{solution}

\begin{solution}{exo:m2:bond-futures:4}
Avec un rendement de 3.8238~\% et un financement à 3.90~\% sur le prix plein coupon de 99\,563\,440~USD, pendant 97
jours, on perd $99\,563\,440 \times 0.000762 \times 97/360 = \text{USD}~20\,451$, si
l'obligation est livrée comme CTD. On achète tout de même la base pour détenir les options du vendeur~:
si les rendements évoluent assez pour que le CTD change, ou si les options de timing et wildcard
rapportent, la base s'élargit et la position longue gagne~; et si l'obligation devient
special en repo, son financement coûte moins que 3.90~\%.
\end{solution}

\begin{solution}{exo:m2:bond-futures:5}
Le future évolue comme le CTD divisé par son facteur, donc son DV01 par contrat
égal au DV01 du CTD par 100\,000~USD, divisé par le facteur. Nombre de
$= \frac{100\,000\,000/100\,000 \times \mathrm{DV01}_{\text{CTD}}}{\mathrm{DV01}_{\text{CTD}}/\mathrm{CF}} = 1\,000 \times 0.8870 = 887$
contrats. Le facteur est le nombre de futures «~par obligation~».
\end{solution}

\begin{solution}{exo:m2:bond-futures:6}
Le facteur valorise chaque obligation à 6~\%. En dessous de 6~\%, toutes les obligations valent plus que
leur valeur ajustée par le facteur, et plus la duration est longue, plus elles ont
gagné à mesure que les rendements baissaient depuis 6~\%~: les obligations longues sont les plus chères à livrer, les
plus courtes les moins chères. Quand les rendements remontent vers 6~\%, les obligations longues perdent davantage,
l'écart se réduit et, à un certain rendement (inférieur à 6~\% ici, car la courbe et les
coupons diffèrent), l'obligation la plus longue devient la moins chère.
\end{solution}

\begin{solution}{exo:m2:bond-futures:7}
L'écart-type sur 97 jours vaut $80 \times \sqrt{97/365} = 41$ points de base.
L'option de changement vaut 0.001 d'un 64e~: le changement est à 3.6
écarts-types. À trois fois la volatilité, soit 124 points de base, elle vaut
3.07 64e. L'\omterm{def:m2:bond-futures:options}{option de qualité} ne vaut rien tant que le changement n'est pas à portée,
puis croît vite~: c'est une option hors de la monnaie sur le niveau de
rendements.
\end{solution}

\begin{solution}{exo:m2:bond-futures:8}
Le future est coté pour 100 d'une obligation notionnelle à 6~\%~; les obligations livrables rapportent
environ 4~\% et valent environ 100 à des rendements de 4~\%. Une obligation à coupon de 6~\% de même
maturité vaudrait environ 112 à ces rendements. Les facteurs (de 0.88 à 0.92)
convertissent~: $111.77 \times 0.887 = 99.14$, à un centième près du 99.13 du CTD. La comparaison doit être
faite après multiplication par le facteur, où l'écart ne dépasse pas quelques centièmes
de point, et non 11~\%.
\end{solution}

\begin{solution}{pb:m2:bond-futures:1}
\textbf{1.} Du 1\up{er} décembre 2026~: août-33 à 6 ans 8 mois, nov.-33 à 6 ans 11
mois, févr.-34 à 7 ans 2 mois, mai-34 à 7 ans 5 mois, août-34 à 7 ans 8
mois~: toutes ont au moins 6 ans 6 mois et moins de 8 ans, et toutes ont été émises comme
obligations d'au plus dix ans~: toutes sont livrables.
\textbf{2.} août-34~: 0.8806. CTD (août-33)~: 0.8870.
\textbf{3.} \omterm{def:m2:bond-futures:irr}{Base brute} $-0.004$~; carry $-0.025$~; \omterm{def:m2:bond-futures:irr}{base nette} $-0.004 + 0.025 = 0.020$.
\textbf{4.} 3.82~\%, huit points de base sous le repo.
\textbf{5.} Le coupon encaissé sur les 97 jours, 1.021 pour 100, est inférieur aux
intérêts du repo sur son prix plein coupon, 1.046~: un coupon de 3.875~\% par an sur le
nominal s'accumule plus lentement que les 3.90~\% facturés sur une année de 360 jours.
\textbf{6.} nov.-33 0.270, févr.-34 0.527, mai-34 0.742, août-34 1.171 points au-dessus du prix
du CTD, 111.793.
\textbf{7.} $+147$ points de base~; l'obligation à 4~\% d'août 2034, la plus longue,
prend le relais.
\textbf{8.} $4.02\% + 1.47\% = 5.49\%$ (l'août-34 rapporte alors 5.55~\%).
\textbf{9.} Les obligations ne rapportent pas toutes le même taux~: les plus longues rapportent davantage, puisque
la courbe est croissante, et elles diffèrent par leur coupon. Le changement a lieu là où le
supplément de rendement de l'obligation la plus longue compense le biais du facteur, en dessous des 6~\% où
une courbe plate le placerait.
\textbf{10.} Le DV01 du future suit celui du nouveau CTD par facteur~: environ
586 à 658 pour 100 de notionnel en futures (en unités de $10^{-4}$), soit 12~\% de plus. Un
hedger qui détient un nombre fixe de contrats se retrouve soudain surcouvert.
\textbf{11.} $80\sqrt{97/365} = 41$ points de base.
\textbf{12.} $147/41 = 3.6$ écarts-types.
\textbf{13.} 0.001 d'un 64e~: sans valeur.
\textbf{14.} 3.07 64e.
\textbf{15.} Les options de timing et wildcard, l'option de fin de mois, et
l'incertitude du marché sur le \omterm{def:m2:repo-and-specials:rate}{taux de repo} auquel le CTD sera financé~: le CTD peut
devenir special et son forward être inférieur au forward sur collatéral général.
\textbf{16.} Acheteuse de la base, elle est longue de l'obligation et courte du future, ce qui
revient à être vendeuse des options que détient le vendeur de futures~: elle les détient. Les options gagnent avec
la volatilité~: un grand mouvement, dans un sens ou dans l'autre, peut rendre une autre
obligation moins chère, et le prix de son obligation longue monte alors par rapport au future.
\textbf{17.} Parce que la duration du future saute lorsque le CTD change~: un ratio de
couverture calculé à partir de l'ancien CTD est faux après le changement, justement au moment où les rendements
ont beaucoup bougé.
\textbf{18.} Une offre plus abondante du CTD~: il devient moins cher sur le marché au comptant et en repo
(moins special), sa \omterm{def:m2:bond-futures:irr}{base nette} baisse, et il devient encore plus nettement le
\omterm{def:m2:bond-futures:ctd}{moins cher à livrer}.
\textbf{19.} Résultat nommé~: \emph{le point de changement} est une hausse parallèle de 147
points de base (rendement du CTD de 5.49~\%), lorsque le 4~\% d'août 2034 remplace le
3.875~\% d'août 2033~; l'option de changement vaut 0.001 d'un 64e à 80 points
de base de volatilité annuelle, et 3.07 64e à 240.
\textbf{20.} L'obligation la moins chère de son \omterm{def:m2:bond-futures:basket}{panier livrable}, divisée par son facteur, diminuée de la
valeur des options du vendeur.
\end{solution}

\begin{solution}{iq:m2:bond-futures:1}
Le contrat permet de livrer plusieurs obligations, de coupons et d'échéances
différents. Le facteur, à savoir le prix de l'obligation pour un rendement de 6~\%, ajuste
le \omterm{def:m2:bond-futures:cf}{prix de facturation} de sorte que chaque obligation soit livrée à un prix proportionnel à sa
valeur si les rendements valaient 6~\%. Sans lui, le vendeur livrerait toujours l'obligation
la moins chère, et le contrat porterait sur une seule émission, exposée aux squeezes.
\iqlookfor{ce qu'est le facteur (un prix à 6~\%) et pourquoi un panier en a besoin.}
\end{solution}

\begin{solution}{iq:m2:bond-futures:2}
Pour chaque livrable, calculer le \omterm{def:m2:bond-futures:irr}{taux de repo implicite} à partir de son prix, du prix à terme,
du facteur, du coupon couru et des coupons jusqu'à la livraison~; le plus élevé est
le CTD. De façon équivalente, c'est la \omterm{def:m2:bond-futures:irr}{base nette} la plus faible. À la livraison, le prix le plus bas
divisé par le facteur.
\iqlookfor{\omterm{def:m2:bond-futures:irr}{taux de repo implicite}, et non le prix le plus bas ni la \omterm{def:m2:bond-futures:irr}{base brute} la plus faible.}
\end{solution}

\begin{solution}{iq:m2:bond-futures:3}
Le taux de financement auquel l'achat de l'obligation, son financement et sa livraison
à terme se compensent exactement. Si le \omterm{def:m2:bond-futures:irr}{taux de repo implicite} du CTD dépasse le \omterm{def:m2:repo-and-specials:rate}{taux de repo}
effectif, acheter la base puis livrer assure un profit certain~: le contrat à terme est cher.
En pratique, le \omterm{def:m2:bond-futures:irr}{taux de repo implicite} du CTD se situe un peu sous le \omterm{def:m2:repo-and-specials:rate}{taux de repo}, l'écart
rémunérant les options de livraison du vendeur.
\iqlookfor{l'interprétation en termes d'arbitrage et la raison pour laquelle l'état normal
se situe sous le \omterm{def:m2:repo-and-specials:rate}{taux de repo}.}
\end{solution}

\begin{solution}{iq:m2:bond-futures:4}
Acheter le CTD, le financer en repo, vendre des contrats à terme au ratio du facteur~;
gagner le \omterm{def:m2:bond-futures:irr}{taux de repo implicite} face au \omterm{def:m2:repo-and-specials:rate}{taux de repo} payé, soit quelques points de base,
démultiplié plusieurs fois par le repo. Ce qui peut mal tourner~: les \omterm{def:m2:repo-and-specials:rate}{taux de repo} s'envolent ou les prêteurs relèvent les décotes,
si bien que le financement coûte plus cher ou doit être réduit~; les appels de marge sur la position vendeuse lors d'une hausse
doivent être honorés en liquidités alors que le gain sur les obligations n'est pas encore réalisé~; la base s'élargit dans une \omterm{def:m2:the-treasury-market:dash}{ruée
vers la liquidité} (dash for cash)~; et le dénouement simultané de nombreux fonds à effet de levier peut faire bouger
le marché où ils se dénouent.
\iqlookfor{le mécanisme du financement et des appels de marge, et le risque de liquidité
de financement, et non seulement l'arbitrage.}
\end{solution}

\begin{solution}{iq:m2:bond-futures:5}
Au premier ordre, le DV01 du CTD pour 100 de nominal, divisé par son \omterm{def:m2:bond-futures:cf}{facteur de concordance},
multiplié par 1\,000 pour un contrat de 100\,000~USD. Mieux~: réévaluer le contrat à terme à partir
de tout le panier sous des rendements choqués (le minimum des prix à terme par facteur, plus un
modèle d'options), ce qui capte la variation du CTD et l'optionnalité~; les deux ne concordent
que loin d'un point de changement.
\iqlookfor{le DV01 rapporté au facteur, et pourquoi le point de changement le rend dépendant de l'état.}
\end{solution}

\begin{solution}{iq:m2:bond-futures:6}
L'\omterm{def:m2:bond-futures:options}{option de qualité} (quelle obligation livrer), l'option de timing (quel jour du mois de livraison),
l'\omterm{def:m2:bond-futures:options}{option wildcard} (livrer après la fixation du prix à terme du jour, ou après la clôture des échanges à un \omterm{def:m2:bond-futures:cf}{prix de facturation} fixe),
et l'option de fin de mois. L'\omterm{def:m2:bond-futures:options}{option de qualité} est généralement la plus importante~: on l'évalue en simulant la
courbe des rendements jusqu'à la livraison (un modèle à facteurs de mouvements parallèles, de pente et de courbure),
en calculant le prix de chaque obligation par facteur, puis en prenant le minimum espéré
face à celui du CTD actuel~; la volatilité et la distance au point de changement
en sont les moteurs.
\iqlookfor{nommer les options et une méthode concrète pour l'\omterm{def:m2:bond-futures:options}{option de qualité}.}
\end{solution}
